% =============================================================================
% Performance profiles for the second-order radius rules on CUTEst.
%
% Standalone. Own preamble, no bibliography, no \cite. The macro block mirrors
% the one near the top of _Thesis_FINAL/Article3/Survey-part3-v1.tex, through
% \providecommand, so this file can later be folded into that document.
% =============================================================================
\documentclass[11pt,a4paper]{article}

\usepackage[utf8]{inputenc}
\usepackage[british]{babel}
\usepackage{amsmath, amssymb, amsthm}
\usepackage{booktabs}
\usepackage{graphicx}
\usepackage{geometry}
\usepackage{float}
\usepackage{enumitem}
\usepackage{subcaption}
\usepackage[hidelinks]{hyperref}
\geometry{margin=2.5cm}

% ---- macros mirrored from Survey-part3-v1.tex -------------------------------
\providecommand{\defeq}{\mathrel{:=}}
\providecommand{\kbar}{\overline{\kappa}}
\providecommand{\calO}{{\mathcal{O}}}
\providecommand{\NN}{{\mathbb{N}}}
\providecommand{\abs}[1]{\left \vert #1 \right \vert}
\providecommand{\norm}[1]{\left \lVert #1 \right \rVert}
\providecommand*\mykappa[1]{ \kappa_{\scriptscriptstyle \text{#1}} }
\providecommand*\myeta[1]{ \eta_{\scriptscriptstyle \text{#1}} }
\providecommand{\Rdelta}{\textsf{R-delta}}
\providecommand{\Rstep}{\textsf{R-step}}
\providecommand{\Rdfo}{\textsf{R-DFO}}
\providecommand{\Rgrad}{\textsf{R-grad}}
\providecommand{\Rrtr}{\textsf{R-RTR}}
\providecommand{\Rgrtr}{\textsf{R-gRTR}}
\providecommand{\RAdaptiveGrad}{\textsf{R-adaptive-grad}}
\providecommand{\RAdaptstep}{\textsf{R-adaptive-step}}
\providecommand{\gone}{\gamma_1}
\providecommand{\gtwo}{\gamma_2}
\providecommand{\gthree}{\gamma_3}

\newtheorem{proposition}{Proposition}
\theoremstyle{remark}
\newtheorem{remark}{Remark}

\title{Performance profiles for the second-order radius rules\\ on CUTEst}
\author{}
\date{}

\begin{document}
\maketitle

\begin{abstract}
PLACEHOLDER-ABSTRACT
\end{abstract}

% =============================================================================
\section{What is measured}\label{sec: measured}
% =============================================================================

Three performance profiles, in this order of importance.

\begin{description}[leftmargin=2.2em, style=nextline]
  \item[P1] Among the second-order radius rules, on the second-order task. Only
    the radius rule differs across the lines, so this profile needs no further
    argument.
  \item[P2] Second-order against first-order, scored on the \emph{first-order}
    task. This is the overhead of carrying the curvature machinery when a
    first-order point is all that was wanted.
  \item[P3] Second-order against first-order, scored on the \emph{second-order}
    task. This is the payoff.
\end{description}

P2 and P3 together are how the two families compare. Neither alone is an answer.
A report quoting only P2 says the second-order arm is expensive, a report quoting
only P3 says it is reliable, and each has described half of one trade.

\subsection{The rule that makes the comparison valid}\label{sec: scoring}

Every run is scored afterwards, from the point it returned, against one criterion
applied by the same code to every run. No profile in this report is drawn on the
solver's own status.

The reason is that the two families stop on different tests. A first-order radius
rule stops when the gradient is small. A second-order rule keeps going until it
also has curvature. Profiling the two on their own stopping tests compares the
cost of an easy task with the cost of a hard one, and the first-order rules win by
construction.

Two tolerances are fixed once,
\begin{equation}
  \varepsilon_g = 10^{-5}, \qquad \varepsilon_H = 10^{-6} .
  \label{eqn: tolerances}
\end{equation}
Both are the package's shared settings, $\varepsilon_g$ the \texttt{tol} of
\texttt{SOLVER\_PARAMS} and $\varepsilon_H$ the \texttt{tol\_H} of
\texttt{SECOND\_ORDER\_PARAMS}. The two criteria are
\begin{align}
  \text{first-order} &: \quad \norm{g} \leq \varepsilon_g ,
  \label{eqn: first order criterion} \\
  \text{second-order} &: \quad \norm{g} \leq \varepsilon_g
     \ \text{ and } \ \lambda_{\min} \geq -\varepsilon_H .
  \label{eqn: second order criterion}
\end{align}
Both are evaluated by \texttt{second\_order\_status} of the package, which is the
same function the solver's own stopping test uses, called here on the returned
point of every run whatever that run's status was.

% =============================================================================
\section{Three things established before the runs}\label{sec: established}
% =============================================================================

\subsection{The subproblem solver}\label{sec: subsolver}

Part~II asks the step for a fraction of the decrease available along the leftmost
eigenvector. Assumption \texttt{as: Eigenpoint decrease} requires
\begin{equation}
  m_k(0) - m_k(s_k) \;\geq\; \tfrac12 \mykappa{fed} \abs{\lambda_k^-} \Delta_k^2
  \qquad\text{whenever } \lambda_k^- < 0 ,
  \label{eqn: eigenpoint decrease}
\end{equation}
and \texttt{as: max Cauchy/eigen decrease} asks for the better of that and the
Cauchy decrease. Theorem \texttt{thm: Eigenpoint decrease} of Part~II supplies
\eqref{eqn: eigenpoint decrease} for the eigenpoint $s_k^E = \Delta_k u_k$ of
\texttt{eqn: eigenpoint}, with $u_k$ any direction satisfying
\texttt{eqn: sufficient negative curvature direction}.

Three subproblem solvers of the package were examined against that condition.

\begin{description}[leftmargin=2.2em, style=nextline]
  \item[\texttt{ExactMS}] Solves the trust-region subproblem exactly by the
    Mor\'e--Sorensen characterisation, hard case included, so it dominates both
    the Cauchy point and the eigenpoint and satisfies
    \eqref{eqn: eigenpoint decrease} with $\mykappa{fed} = 1$. It needs a dense
    eigendecomposition and refuses above its \texttt{nmax}, which is $200$.
  \item[\texttt{EigenPoint(SteihaugCG())}] Computes $\pm\Delta v_{\min}$ and
    returns whichever of that and the inner step decreases the model more, so it
    satisfies \eqref{eqn: eigenpoint decrease} with $\mykappa{fed} = \mykappa{snc}$,
    the quality of the eigenvector. That is $1$ in the dense branch and
    Lanczos-limited above it.
  \item[\texttt{SteihaugCG}] Terminates on the first direction of negative
    curvature the conjugate-gradient recurrence reaches and runs to the boundary
    along it. The decrease is real and the direction is whichever the recurrence
    happened to find, so the guaranteed decrease carries no
    $\abs{\lambda_k^-}\Delta_k^2$ term. It does not deliver the eigenpoint
    decrease.
\end{description}

The second-order arm therefore uses \texttt{ExactMS} at $n \leq 200$ and
\texttt{EigenPoint(SteihaugCG())} above. The switch is on the problem, not on the
column, so every line of a profile sees the same subproblem solver on a given
problem and the per-problem normalisation is unaffected. The first-order arm uses
the package's default, \texttt{SteihaugCG}. Every caption names the pair.

\subsection{The criticality measure}\label{sec: tau}

Part~II defines the second-order criticality measure at
\texttt{eqn:tau definition},
\begin{equation}
  \tau_k \defeq \max\{\norm{g_k},\, -\lambda_k^-\} \;\geq\; 0 ,
  \label{eqn: tau}
\end{equation}
with $\lambda_k^-$ the leftmost eigenvalue of the \emph{model} Hessian $H_k$.
Part~II states that $\tau_k$ is a property of the pair $(x_k, H_k)$ and not of
$x_k$ alone, for that reason.

The package computes $\tau$ in \texttt{tau\_criticality} as
$\max\{\norm{g}, -\min(\lambda_{\min}, 0)\}$. The inner clamp is a no-op inside
the maximum, since $\norm{g} \geq 0$. The two forms were compared on a grid of
twenty pairs $(\norm{g}, \lambda)$ covering $\lambda < 0$, $\lambda = 0$ and
$\lambda > 0$, and the worst absolute difference is $0$. So the package
implements \eqref{eqn: tau} exactly and no choice arises.

The open note at \texttt{app: measure} of the numerical study says that $\tau_k$
is the leftmost model Hessian eigenvalue in Part~I and the whole second-order
measure in Part~II, and that Part~III must pick one. \textbf{The premise is
inaccurate and there is nothing to pick.} Part~I defines no second-order
criticality measure. Its criticality-anchored section carries a generic measure
$\chi_k$ with $\chi_k = \norm{g_k}$ as its stated example, at
\texttt{eqn:criticality step}, and its concluding section says that second-order
stationarity is outside the scope of its framework. The string \texttt{eigen} does
not occur anywhere in the Part~I source, and its four occurrences of
\texttt{lambda\_} are the rates $\lambda_1$ and $\lambda_2$ of the adaptive
update. The note should be resolved by recording that $\tau_k$ is
\eqref{eqn: tau} of Part~II, which is what the package already implements.

\subsection{Per-iteration costs, and the route taken}\label{sec: costs}

Post hoc scoring needs the cost a run had spent when it \emph{first} met the
criterion, not its total at the end. The package's trace does not carry that. The
trajectories it records are the radius, the gradient norm, the objective, the
ratio, the step norm, the activity and acceptance flags, $\tau$, the curvature
and a further nine diagnostics. Not one of them is an evaluation count. The benchmark record adds three totals, read off the model after the solve,
and no per-iteration counts either.

The fallback route was taken. Each configuration is run twice, with the tightened
criterion as its \emph{actual} stopping rule.

\begin{description}[leftmargin=2.2em, style=nextline]
  \item[Pass g] halts at the first iterate with $\norm{g_k} \leq \varepsilon_g$.
    Its cost is what P2 scores.
  \item[Pass h] halts at the first iterate that also has
    $\lambda_{\min} \geq -\varepsilon_H$. Its cost is what P1 and P3 score.
\end{description}

This is more runs and it is exact. Both passes keep their arm's own
\texttt{tol\_H}, which in this package controls the curvature estimate as well as
the stopping test. Switching it off in pass~g would also switch off the cost that
pass exists to measure, so pass~g halts through a callback instead.

One arithmetic consequence of the package's design bears on
Section~\ref{sec: cost axes}. With \texttt{ExactHessian} at $n \leq 200$ the
curvature estimate goes through \texttt{dense\_hessian}, hence through
\texttt{NLPModels.hess}, and is paid in \emph{Hessian evaluations}. Above $200$ it
goes through Lanczos on \texttt{hessian\_op} and is paid in Hessian-vector
products. A profile on Hessian-vector products alone would therefore show the
second-order arm as free on the small problems, which is the opposite of true.
Both axes are reported.

% =============================================================================
\section{The configurations}\label{sec: configs}
% =============================================================================

Six radius rules, in two groups of three.

\begin{table}[H]
  \centering
  \caption{The roster. The second column is \texttt{is\_criticality\_anchored} and
    the third \texttt{asymptotic\_regime}, both read from the package rather than
    assumed. $\omega_k$ is the measure the radius is anchored to in the
    second-order arm.}
  \label{tab: roster}
  \begin{tabular}{llll}
    \toprule
    rule & reads a measure & regime & $\omega_k$ in the second-order arm \\
    \midrule
    $\Rdelta$ & no  & \texttt{bounded\_below} & not read \\
    $\Rstep$  & no  & \texttt{step\_summable} & not read \\
    $\Rrtr$   & no  & \texttt{step\_summable} & not read \\
    $\Rdfo$   & yes & \texttt{vanishing}      & $\tau_k$ \\
    $\Rgrad$  & yes & \texttt{vanishing}      & $\tau_k$ \\
    $\Rgrtr$  & yes & \texttt{vanishing}      & $\tau_k$ \\
    \bottomrule
  \end{tabular}
\end{table}

The distinction in Table~\ref{tab: roster} governs how every caption in this
report must be read. $\Rdfo$, $\Rgrad$ and $\Rgrtr$ read a criticality measure, so
replacing $\norm{g_k}$ by $\tau_k$ changes the radius rule itself. $\Rdelta$,
$\Rstep$ and $\Rrtr$ do not read one, so for them the second-order arm is
\emph{the same radius rule under a different stopping test and a different
subproblem solver}. Both groups belong in P1, and a caption that does not say
which is which claims a distinction the figure does not have.

The two arms differ in exactly three things.

\begin{table}[H]
  \centering
  \caption{The two arms. Everything not listed is identical, including the model
    Hessian, the acceptance and scaling thresholds, the contraction and expansion
    factors, the initial radius, the iteration cap and the two tolerances of
    \eqref{eqn: tolerances}.}
  \label{tab: arms}
  \begin{tabular}{lll}
    \toprule
     & first-order arm & second-order arm \\
    \midrule
    second-order stopping test & off & on, $\varepsilon_H = 10^{-6}$ \\
    $\omega_k$, where read     & $\norm{g_k}$ & $\tau_k$ \\
    subproblem solver          & \texttt{SteihaugCG}
                               & \texttt{ExactMS} ($n \leq 200$), \\
                               & & \texttt{EigenPoint(SteihaugCG())} ($n > 200$) \\
    \bottomrule
  \end{tabular}
\end{table}

The shared settings are those of \texttt{SOLVER\_PARAMS}, unchanged:
$\eta = \eta_1 = 0.1$, $\eta_2 = 0.9$, $\gone = 0.25$, $\gtwo = 0.5$,
$\gthree = 2$, $\Delta_0 = 1$, $\Delta_{\min} = 0$, and at most $10^4$ iterations.
$\Rdfo$ runs at $\zeta = 100$ and $\Rgrad$ and $\Rgrtr$ at $\mu_0 = 1$ with no
cap. The two retrospective rules run at their own $\tilde\eta_1 = 0.05$ and
$\tilde\eta_2 = 0.9$. Per-rule tuning would measure tuning effort rather than
algorithmic merit, which is why the settings are shared.

The model Hessian is \texttt{ExactHessian} in both arms. That choice is what makes
the criterion \eqref{eqn: second order criterion} a statement about $f$. It was
checked rather than assumed. On the problems tested, the model's dense Hessian
agrees with \texttt{NLPModels.hess} entry by entry, and the package's
$\lambda_{\min}$ agrees with a direct dense eigendecomposition, both to $0$ in
floating point.

\subsection{\texttt{LBFGSModel} gets no line}\label{sec: lbfgs}

\texttt{LBFGSModel} enforces $B \succ 0$, so $\lambda_{\min} > 0$ at every
iteration, so $\tau_k = \norm{g_k}$ identically and the second-order variant of a
rule \emph{is} its first-order variant. It therefore has no line of its own in P1.

This was verified on iterate sequences and not asserted. For each of the three
criticality-anchored rules, on two problems, the wrapped and unwrapped rules were
run over \texttt{LBFGSModel} and their iteration counts, gradient trajectories and
radius trajectories compared. All six pairs are identical, entry by entry. There
is no difference to report as a finding.

% =============================================================================
\section{The problem set}\label{sec: problems}
% =============================================================================

CUTEst, unconstrained, with all variables free, $2 \leq n \leq 1000$ and no
constraints. CUTEst version \texttt{1.4.0}. The query returns \textbf{196}
problems, of which \textbf{185} have $n \leq 200$ and \textbf{11} have $n > 200$.
The realised list of names and dimensions is written beside the results, in
\texttt{exp17\_problems.txt}.

The fallback was disabled deliberately. \texttt{default\_problems()} of the
benchmark returns the analytic set when the CUTEst query comes back empty, so a
run in which CUTEst failed to load produces a table that reads as a CUTEst
benchmark and is not one. The experiment checks \texttt{HAS\_CUTEST}, checks that
the query returned something, and raises rather than falling back.

P3 runs on the $185$ problems with $n \leq 200$ and no others. Above that
dimension the package estimates $\lambda_{\min}$ by Lanczos, a Ritz value
satisfies $\lambda_{\text{Ritz}} \geq \lambda_{\min}$, and an under-resolved run
therefore \emph{understates} the negative curvature and makes the second-order
test optimistic. Scoring P3 on such an estimate would rank the runs by the
machinery the profile exists to compare. On the $185$ problems the estimate is a
dense symmetric eigendecomposition of $\nabla^2 f(x)$ and is exact.

PLACEHOLDER-RESULTS

\end{document}
